\subsection{交换条件}

``交换条件''是关于 (W, S) 的如下断言：

\begin{enumerate}
\def\labelenumi{(\Alph{enumi})}
\setcounter{enumi}{4}
\tightlist
\item
  设 \(w \in W\) 和 \(s \in S\) 满足 \(l(sw) \leqslant l(w)\)。对于任意约化分解 \((s_1, \ldots, s_q)\) 的 \(w\)，存在整数 \(j\)，满足 \(1 \leqslant j \leqslant q\)，并且
\end{enumerate}

\[
s s _ {1} \dots s _ {j - 1} = s _ {1} \dots s _ {j - 1} s _ {j}.\tag{16}
\]

在本号中，我们假定 \((W,S)\) 满足 (E)。由引理 3，当 \((W,S)\) 是 Coxeter 系统时，这一条件成立。因此，本号的结果适用于 Coxeter 系统。

命题 4. 设 \(s \in S\)、\(w \in W\)，且 \(s = (s_1, \ldots, s_q)\) 是 \(w\) 的一个约化分解。则下列两种情形必有一种成立：

\begin{enumerate}
\def\labelenumi{\alph{enumi})}
\item
  \(l(sw) = l(w) + 1\)，且 \((s, s_1, \ldots, s_q)\) 是 sw 的一个约化分解。
\item
  \(l(sw) = l(w) - 1\)，并且存在整数 \(j\)，满足 \(1 \leqslant j \leqslant q\)，\((s_1, \ldots, s_{j-1}, s_{j+1}, \ldots, s_q)\) 是 \(sw\) 的一个约化分解，而 \((s, s_1, \ldots, s_{j-1}, s_{j+1}, \ldots, s_q)\) 是 \(w\) 的一个约化分解。
\end{enumerate}

置 \(w' = sw\)。由第 1 号的公式 (3)，有

\[
\left| l (w) - l \left(w ^ {\prime}\right) \right| \leqslant l (s) = 1.
\]

我们区分两种情形：

\begin{enumerate}
\def\labelenumi{\alph{enumi})}
\tightlist
\item
  \(l(w') > l(w)\)。于是 \(l(w') = q + 1\) 且 \(w' = ss_1 \ldots s_q\)，所以
\end{enumerate}

\[
\left(s, s _ {1}, \ldots , s _ {q}\right)
\]

是 \(w'\) 的一个约化分解。

\begin{enumerate}
\def\labelenumi{\alph{enumi})}
\setcounter{enumi}{1}
\tightlist
\item
  \(l(w') \leqslant l(w)\)。由 (E)，存在整数 \(j\)，满足 \(1 \leqslant j \leqslant q\)，且 (16) 成立。于是 \(w = ss_1 \ldots s_{j-1}s_{j+1} \ldots s_q\)，从而
\end{enumerate}

\[
w ^ {\prime} = s _ {1} \dots s _ {j - 1} s _ {j + 1} \dots s _ {q}.
\]

由于 \(q - 1 \leqslant l(w') \leqslant q\)，可知 \(l(w') = q - 1\)，且 \((s_1, \ldots, s_{j-1}, s_{j+1}, \ldots, s_q)\) 是 \(w'\) 的一个约化分解。

引理 4. 设 \(w \in W\) 的长度为 \(q \geqslant 1\)，令 \(D\) 为 \(w\) 的约化分解集合，并令 \(F\) 为从 \(D\) 到集合 \(E\) 的映射。假设 \(F(s) = F(s')\)，其中 \(s = (s_1, \ldots, s_q)\)、\(s' = (s_1', \ldots, s_q')\) 是 \(D\) 中满足下列某一假设的元素：

\[
a) s _ {1} = s _ {1} ^ {\prime} \text {   or   } s _ {q} = s _ {q} ^ {\prime}.
\]

\begin{enumerate}
\def\labelenumi{\alph{enumi})}
\setcounter{enumi}{1}
\tightlist
\item
  存在 \(s\) 和 \(s'\)，属于 \(S\)，使得 \(s_j = s'_k = s\) 且 \(s_k = s'_j = s'\)，其中 \(j\) 为奇数且 \(k\) 为偶数。
\end{enumerate}

则 F 为常值映射。

\begin{enumerate}
\def\labelenumi{\Alph{enumi})}
\item
  设 \(\mathbf{s}, \mathbf{s}' \in D\)，并置 \(\mathbf{t} = (s_1', s_1, \ldots, s_{q-1})\)。我们将证明：若 \(\mathrm{F}(\mathbf{s}) \neq \mathrm{F}(\mathbf{s}')\)，则 \(\mathbf{t} \in D\) 且 \(\mathrm{F}(\mathbf{t}) \neq \mathrm{F}(\mathbf{s})\)。事实上，\(w = s_1' \ldots s_q'\)，因而 \(s_1' w = s_2' \ldots s_q'\) 的长度为 \(< q\)。由命题 4，存在整数 \(j\)，满足 \(1 \leqslant j \leqslant q\)，使序列 \(\mathbf{u} = (s_1', s_1, \ldots, s_{j-1}, s_{j+1}, \ldots, s_q)\) 属于 D。有 \(\mathrm{F}(\mathbf{u}) = \mathrm{F}(\mathbf{s}')\)，根据条件 \(a\)；若 \(j \neq q\)，同理有 \(\mathrm{F}(\mathbf{s}) = \mathrm{F}(\mathbf{u})\)，从而 \(\mathrm{F}(\mathbf{s}) = \mathrm{F}(\mathbf{s}')\)，这与假设矛盾。因此 \(j = q\)，于是 \(\mathbf{t} = \mathbf{u} \in D\)，且 \(\mathrm{F}(\mathbf{t}) = \mathrm{F}(\mathbf{s}') \neq \mathrm{F}(\mathbf{s})\)。
\item
  设 \(\mathbf{s}\) 和 \(\mathbf{s}'\) 属于 D。对于任意整数 \(j\)，满足 \(0 \leqslant j \leqslant q + 1\)，定义一个序列 \(\mathbf{s}_j\)，由 S 中的 \(q\) 个元素组成，如下：
\end{enumerate}

\[
\begin{array}{r l} \mathbf {s} _ {0} & = (s _ {1} ^ {\prime}, \dots , s _ {q} ^ {\prime}) \\ \mathbf {s} _ {1} & = (s _ {1}, \dots , s _ {q}) \\ \mathbf {s} _ {q + 1 - k} & = (s _ {1}, s _ {1} ^ {\prime}, \dots , s _ {1}, s _ {1} ^ {\prime}, s _ {1}, s _ {2}, \dots , s _ {k}) \\ & \text {for } q - k \text { even and } 0 \leqslant k \leqslant q \end{array}\tag{17}
\]

\[
\mathbf {s} _ {q + 1 - k} = \left(s _ {1} ^ {\prime}, s _ {1}, \ldots , s _ {1}, s _ {1} ^ {\prime}, s _ {1}, s _ {2}, \ldots , s _ {k}\right)
\]

\[
\text { for   } q - k \text {   odd   and   } 0 \leqslant k \leqslant q.
\]

记 \((\mathrm{H}_{j})\) 表示断言 `` \(s_{j} \in D, s_{j+1} \in D\) 且 \(F(s_{j}) \neq F(s_{j+1})\) ''。由 A)，\((\mathrm{H}_{j}) \Longrightarrow (\mathrm{H}_{j+1})\)，且这对 \(0 \leqslant j < q\) 成立；而根据条件 b)，\((\mathrm{H}_{q})\) 不成立。因此，\((\mathrm{H}_{0})\) 不成立。由于 \(s_{0} = s'\) 且 \(s_{1} = s\)，得到 \(F(s) = F(s')\)。

命题 5. 设 M 为一个幺半群（单位元为 1），f 为从 S 到 M 的映射。对于 \(s, s' \in S\)，令 \(m(s, s')\) 为 \(ss'\) 的阶，并置

\[
a (s, s ^ {\prime}) = \left\{ \begin{array}{l l} (f (s) f (s ^ {\prime})) ^ {l} & \text {if } m(s,s^{\prime}) = 2l,\ l \text { finite} \\ (f (s) f (s ^ {\prime})) ^ {l} f (s) & \text {if } m(s,s^{\prime}) = 2l + 1,\ l \text { finite} \\ 1 & \text {if } m(s,s^{\prime}) = \infty . \end{array} \right.\tag{18}
\]

如果 \(a(s, s') = a(s', s)\) 对 \(s \neq s'\) 属于 \(S\) 的情形都成立，则存在一个映射 \(g\)，从 \(W\) 到 \(M\)，使得

\[
g (w) = f \left(s _ {1}\right) \dots f \left(s _ {q}\right)\tag{19}
\]

对所有 \(w \in \mathbf{W}\) 以及任意约化分解 \((s_1, \ldots, s_q)\) 的 \(w\)，

对于任意 \(w \in W\)，令 \(D_w\) 为 \(w\) 的约化分解集合，并令 \(F_w\) 为从 \(D_w\) 到 \(M\) 的映射，由下式定义：

\[
\mathrm{F} _ {w} (s _ {1}, \dots , s _ {q}) = f (s _ {1}) \dots f (s _ {q}).
\]

我们将按 w 的长度归纳，证明 \(F_{w}\) 为常值，这将证明命题 5。由于 \(l(w) = 0, 1\) 的情形是平凡的，我们假定 \(q \geqslant 2\)，并假定对于满足 \(l(w) < q\) 的元素 w，断言已经得到证明。设 w 的长度为 q，且 \(s, s' \in D_w\)；根据引理 4，只需证明在该引理的 a) 和 b) 两种情形下 \(\mathrm{F}_{w}(\mathbf{s}) = \mathrm{F}_{w}(\mathbf{s}')\)。

\begin{enumerate}
\def\labelenumi{\alph{enumi})}
\tightlist
\item
  公式
\end{enumerate}

\[
\mathrm{F} _ {w} \left(s _ {1}, \dots , s _ {q}\right) = f \left(s _ {1}\right) \mathrm{F} _ {w ^ {\prime \prime}} \left(s _ {2}, \dots , s _ {q}\right) = \mathrm{F} _ {w ^ {\prime}} \left(s _ {1}, \dots , s _ {q - 1}\right) f \left(s _ {q}\right)
\]

当 \(w' = s_1 \ldots s_{q-1}\) 且 \(w'' = s_2 \ldots s_q\) 时，上式以及归纳假设表明：\(F_w(\mathbf{s}) = F_w(\mathbf{s}')\)，若 \(s_1 = s_1'\) 或 \(s_q = s_q'\)。

\begin{enumerate}
\def\labelenumi{\alph{enumi})}
\setcounter{enumi}{1}
\tightlist
\item
  假设存在两个元素 \(s\) 和 \(s'\)，属于 \(S\)，使得 \(s_j = s'_k = s\) 且 \(s_k = s'_j = s'\)，其中 \(j\) 为奇数且 \(k\) 为偶数。只需处理 \(s \neq s'\) 的情形。此时，序列 \(s\) 和 \(s'\) 是 \(w\) 的两个不同约化分解，属于由 \(s\) 和 \(s'\) 生成的二面体群。根据第 2 号备注，\(m\)（即 \(ss'\) 的阶）必为有限数；用该备注的记号，\(s = s_m\) 且 \(s' = s'_m\)。因此，\(F_w(s) = a(s, s')\) 且 \(F_w(s') = a(s', s)\)，于是
\end{enumerate}

\[
\mathrm{F} _ {w} (\mathbf {s}) = \mathrm{F} _ {w} (\mathbf {s} ^ {\prime}).
\]

